\section{Zadanie 7}
Powiedzmy, że w urnie znajduje się $N$ kul z czego $K$ jest zielonych a reszta czerwone. Losujemy $n$ kul bez zwracania. 
Niech $X$ będzie liczbą zielonych kul które wylosowaliśmy. Jaki rozkład ma $X$? Innymi słowy, ile wynosi $P(X = k)$?

\subsection*{Rozwiązanie}
$k$ kul zielonych spośród wszystkich kul zielonych można wybrać na $\binom{K}{k}$ sposobów, analogicznie $n-k$ kul czerwonych na $\binom{N-K}{n-k}$ sposobów. Wszystkich sposobów na wybranie $n$ kul ze zbioru $N$ kul jest $\binom{N}{n}$. Wynik wynosi więc:

\[
\frac{\binom{K}{k} \cdot \binom{N-K}{n-k}}{\binom{N}{n}}
\]

\subsection*{Odpowiedź}
\[
\boxed{\frac{\binom{K}{k} \cdot \binom{N-K}{n-k}}{\binom{N}{n}} \text{, jest to rozkład hipergeometryczny}}
\]
